\documentclass[12pt, ukrainian]{article}
\usepackage[a4paper]{geometry}
%\setlength\textwidth{170mm} \setlength\hoffset{-25mm}
%\setlength\textheight{277mm} \setlength\voffset{-38mm}
%\setlength\parindent{0mm}
\geometry{top=1.5cm,bottom=1.5cm,inner=1.5cm,outer=1.5cm}
\pagestyle{empty}
\usepackage[T2A]{fontenc}
\usepackage[utf8]{inputenc}
\usepackage[ukrainian]{babel}
\usepackage{graphicx}
\usepackage{caption}
\DeclareCaptionFormat{nocaption}{}
\captionsetup{format=nocaption,aboveskip=0pt,belowskip=0pt}
\usepackage[Export]{adjustbox}
\adjustboxset{max size={0.9\linewidth}{0.9\paperheight}}
\usepackage{verbatim}
\usepackage{fancyvrb}
\def\code#1{\Verb!#1!}
\begin{document}
{\Large{23.10.2024 Контрольна робота \No 1, варіант  \No \textbf { 158 }\par}}
{
\begin {large}
{\begin {center}\textbf{Завдання 1}\end {center}}

Що буде виведено за виконання фрагменту коду?\par
\begin{center}
	\adjustimage{max size={0.9\linewidth}{0.9\paperheight}}
	{../pict/158_1.png}
\end{center}
\newpage

{\begin {center}\textbf{Завдання 2}\end {center}}

Що буде виведено за виконання фрагменту коду?\par
\begin{center}
	\adjustimage{max size={0.9\linewidth}{0.9\paperheight}}
	{../pict/158_2.png}
\end{center}
\newpage

{\begin {center}\textbf{Завдання 3}\end {center}}

Що буде виведено за виконання фрагменту коду?\par
\begin{center}
	\adjustimage{max size={0.9\linewidth}{0.9\paperheight}}
	{../pict/158_3.png}
\end{center}
\newpage

{\begin {center}\textbf{Завдання 4}\end {center}}

Що буде виведено за виконання програми?\par
\begin{center}
	\adjustimage{max size={0.9\linewidth}{0.9\paperheight}}
	{../pict/158_4.png}
\end{center}
\newpage

{\begin {center}\textbf{Завдання 5}\end {center}}

Що буде виведено за виконання фрагменту коду?\par
\begin{center}
	\adjustimage{max size={0.9\linewidth}{0.9\paperheight}}
	{../pict/158_5.png}
\end{center}
\newpage

{\begin {center}\textbf{Завдання 6}\end {center}}

Що буде виведено за виконання фрагменту коду?\par
\begin{center}
	\adjustimage{max size={0.9\linewidth}{0.9\paperheight}}
	{../pict/158_6.png}
\end{center}
\newpage

{\begin {center}\textbf{Завдання 7}\end {center}}

Що буде виведено за виконання програми?\par
\begin{center}
	\adjustimage{max size={0.9\linewidth}{0.9\paperheight}}
	{../pict/158_7.png}
\end{center}
\newpage

{\begin {center}\textbf{Завдання 8}\end {center}}

Що буде виведено за виконання програми?\par
\begin{center}
	\adjustimage{max size={0.9\linewidth}{0.9\paperheight}}
	{../pict/158_8.png}
\end{center}
\newpage

{\begin {center}\textbf{Завдання 9}\end {center}}

Виберіть версію функції \code{f}, за якої за виконання коду
\begin{center}
	\adjustimage{max size={0.9\linewidth}{0.9\paperheight}}
	{../pict/158_90.png}
\end{center}
буде виведено 314

\vspace{5mm}
версія 1 :

\begin{center}
	\adjustimage{max size={0.9\linewidth}{0.9\paperheight}}
	{../pict/158_91.png}
\end{center}
версія 2 :

\begin{center}
	\adjustimage{max size={0.9\linewidth}{0.9\paperheight}}
	{../pict/158_92.png}
\end{center}
версія 3 :

\begin{center}
	\adjustimage{max size={0.9\linewidth}{0.9\paperheight}}
	{../pict/158_93.png}
\end{center}
версія 4 :

\begin{center}
	\adjustimage{max size={0.9\linewidth}{0.9\paperheight}}
	{../pict/158_94.png}
\end{center}
\newpage

{\begin {center}\textbf{Завдання 10}\end {center}}

Що буде виведено за виконання програми?\par
\begin{center}
	\adjustimage{max size={0.9\linewidth}{0.9\paperheight}}
	{../pict/158_10.png}
\end{center}
\newpage

{\begin {center}\textbf{Завдання 11}\end {center}}

Що буде виведено за виконання програми?\par
\begin{center}
	\adjustimage{max size={0.9\linewidth}{0.9\paperheight}}
	{../pict/158_11.png}
\end{center}
\newpage

{\begin {center}\textbf{Завдання 12}\end {center}}


Написати функцiю f, що отримує ціле число і повертає логічну ознаку того, що воно належить вiдрiзку [2021; 2022]. Стиль запису умови також оцінюється.\par
\newpage

\end {large}}
%end
\end{document}

